% GENERATED FILE. Edit canonical lessons, then run npm run generate:books.
% canonical-source-hash: fnv1a64:dca4c6b1460ad67f
% canonical-lessons: ZH-C93-taihaole, ZH-W93-tai, ZH-C93-youdianr, ZH-C93-ganjing, ZH-C93-zang, ZH-C93-practice

\chapter{Too Much}
\label{ch:zh-too-much}
\hlchaptermodality{93}

\begin{chapteropening}
\textbf{By the end of this chapter:} \emph{I can say something is too much, or a bit too much.}

Too much with \zh{太} … \zh{了}, a bit too much with \zh{有点儿}, great with \zh{太好了}, and clean and dirty, built on one new written character: \zh{太}.
\end{chapteropening}

\section[\texorpdfstring{\zh{太好了} (tài hǎo le)}{太好了 (tài hǎo le)}]{\zh{太好了} (tài hǎo le) — great!}

\label{lesson:ZH-C93-taihaole}



\begin{quote}
A new chapter, and it opens with a sound you have made since the chapter on money.

\begin{itemize}
\raggedright
  \item \emph{Recall:} say \zh{我也要} — \emph{wǒ yě yào}, me too
  \item \emph{Recall:} say \emph{nǐ yào biéde ma}, anything else — \textbf{R1}, the chapter you have finished
  \item \emph{Recall:} say \emph{tài guì le}, too expensive — \textbf{R4}, from the chapter on money
\end{itemize}
\end{quote}

\subsection*{You'll want to know: \zh{太好了}}
\textbf{\zh{太好了}} — \emph{tài hǎo le} — \textquotedblleft{}great!; that's wonderful\textquotedblright{}.

\begin{quote}
\textbf{\zh{我们明天去} \zh{太好了}} \emph{wǒmen míngtiān qù, tài hǎo le} — We're going tomorrow? Great!
\end{quote}

The next lesson puts the first character in your hand. It is the \emph{tài} you have said since \emph{tài guì le}.

\begin{cousinweb}
\textbf{\zh{太好了}} (\emph{tài hǎo le}) is built from \zh{太} 'too; extremely', \zh{好} 'good' and \zh{了}. \zh{太} and \zh{好} are inherited from Middle Chinese, \emph{thajH} and \emph{xawX} (Guangyun, in Baxter's transcription), \emph{thajH} glossed 'extreme; great'.
\end{cousinweb}

\begin{grammarlens}[title={What you've built}]
\emph{tài … le} around an adjective says too much: \emph{tài guì le}, too expensive. Around \zh{好} it is not 'too good'. It is the everyday 'great!'
\end{grammarlens}

\subsection*{Your turn}
\begin{itemize}
\raggedright
  \item \emph{Say it:} \emph{tài hǎo le}
  \item \emph{Say it:} \emph{tài hǎo le}, once more, as if you meant it
  \item \emph{Say it:} \zh{我们明天去} \zh{太好了} — \emph{wǒmen míngtiān qù, tài hǎo le} — We're going tomorrow? Great!
  \item \emph{Recall:} say \emph{tài guì le}, then \emph{tài hǎo le}, and hear one complaint and one delight
\end{itemize}

\begin{tcolorbox}[breakable,colback=teal!4,colframe=teal!35!black,title={Before you move on}]
What does \zh{太好了} mean? (\textquotedblleft{}Great!\textquotedblright{}.) Is it 'too good'? (\textbf{No: \emph{tài … le} around \zh{好} is plain delight}.)
\end{tcolorbox}
\section[\texorpdfstring{\zh{太} (tài)}{太 (tài)}]{\zh{太} — \zh{大}, then a dot}

\label{lesson:ZH-W93-tai}



\begin{quote}
The word first, then the character inside it.

\begin{itemize}
\raggedright
  \item \emph{Recall:} say \emph{tài hǎo le} — \textbf{R1}, one lesson back, before it decays
\end{itemize}

\begin{quote}
\zh{大} \zh{天}
\end{quote}

\emph{Read it:} each one aloud, with nothing said first — \emph{dà}, \emph{tiān}
\end{quote}

\subsection*{Script — \zh{大}, then a dot}
\hlblockfigure{\detokenize{figures/ZH-W93-tai-filmstrip.pdf}}{How \zh{太} is written, stroke by stroke}

\begin{quote}
\zh{太}
\end{quote}

\textbf{\zh{大} with one stroke added, and the added stroke comes last.}

\textbf{\zh{大}, three strokes}: the bar, the left-falling stroke through it, the right-falling stroke from the crossing.

\textbf{One dot}, between the two legs.

3 + 1 = four.

\subsection*{Your turn}
\begin{itemize}
\raggedright
  \item \emph{Write it:} \zh{太}, all four, with the model in front of you
  \item \emph{Cover:} write it again from memory, and count the pieces, not the strokes
\end{itemize}

\begin{tcolorbox}[breakable,colback=teal!4,colframe=teal!35!black,title={Before you move on}]
How many strokes in \zh{太}? (\textbf{Four}.) Which character is inside it? (\textbf{\zh{大}}.) Where does the dot go? (\textbf{Between the legs, last}.)

Source: \href{https://raw.githubusercontent.com/chanind/hanzi-writer-data/68d10a4b21150cae5e1ebbd223eed289cf32d90c/data/\%E5\%A4\%AA.json}{Hanzi Writer Data \zh{太}, pinned stroke paths and medians}.
\end{tcolorbox}
\section[\texorpdfstring{\zh{有点儿} (yǒudiǎnr)}{有点儿 (yǒudiǎnr)}]{\zh{有点儿} (yǒudiǎnr) — a bit, and not in a good way}

\label{lesson:ZH-C93-youdianr}



\begin{quote}
Before the new one: read the last lesson's character, then two words this one is made from.

\begin{itemize}
\raggedright
  \item \emph{Read it:} \zh{太好了}, aloud — \emph{tài hǎo le}
  \item \emph{Recall:} say \zh{一点儿} — \emph{yìdiǎnr}, a little
  \item \emph{Recall:} say \emph{lěng}, cold — \textbf{R4}, recognition at distance
\end{itemize}
\end{quote}

\subsection*{You'll want to know: \zh{有点儿}}
\textbf{\zh{有点儿}} — \emph{yǒudiǎnr} — \textquotedblleft{}a bit (of something unwanted); rather\textquotedblright{}.

\begin{quote}
\textbf{\zh{有点儿大}} \emph{yǒudiǎnr dà} — It's a bit big.
\end{quote}

\begin{quote}
\textbf{\zh{太大了}} \emph{tài dà le} — It's too big.
\end{quote}

Two dips in a row: spoken, the first rises, \emph{yóudiǎnr}. It is written \emph{yǒudiǎnr}.

\begin{cousinweb}
\textbf{\zh{有点儿}} (\emph{yǒudiǎnr}) is built from \zh{有} 'to have', \zh{点} 'a dot' and the r ending: having a dot of it. \zh{有} and \zh{点} are inherited from Middle Chinese, \emph{hjuwX} and \emph{temX} (Guangyun, in Baxter's transcription).
\end{cousinweb}

\begin{grammarlens}[title={What you've built}]
\zh{有点儿} + an adjective is a mild complaint: \emph{yǒudiǎnr lěng}, it's a bit cold. \zh{很} is neutral; \zh{太} … \zh{了} is strong. And \zh{一点儿} is a different job: it goes after a verb, \zh{要一点儿水}.
\end{grammarlens}

\subsection*{Your turn}
\begin{itemize}
\raggedright
  \item \emph{Say it:} \emph{yǒudiǎnr} — and hear it as \emph{yóudiǎnr}
  \item \emph{Say it:} \zh{有点儿大} — \emph{yǒudiǎnr dà} — It's a bit big.
  \item \emph{Say it:} \zh{太大了} — \emph{tài dà le} — It's too big.
  \item \emph{Say it:} \emph{yǒudiǎnr lěng} — It's a bit cold.
  \item \emph{Recall:} say \emph{yào yìdiǎnr shuǐ}, then \emph{yǒudiǎnr dà}, and hear the two jobs
\end{itemize}

\begin{tcolorbox}[breakable,colback=teal!4,colframe=teal!35!black,title={Before you move on}]
What does \zh{有点儿} say about an adjective? (\textbf{A bit, and not welcome}.) Which is stronger, \zh{有点儿大} or \zh{太大了}? (\textbf{\zh{太大了}}.)
\end{tcolorbox}
\section[gānjìng]{gānjìng — clean}

\label{lesson:ZH-C93-ganjing}



\begin{quote}
Before the new one: last lesson's word, then a room, then \zh{很}.

\begin{itemize}
\raggedright
  \item \emph{Recall:} say \zh{有点儿大} — \emph{yǒudiǎnr dà}
  \item \emph{Recall:} say \emph{fángjiān}, a room — \textbf{R2}
  \item \emph{Recall:} say \zh{很好} — \emph{hěn hǎo}, fine
\end{itemize}
\end{quote}

\subsection*{You'll want to know: gānjìng}
\emph{gānjìng} — \textquotedblleft{}clean\textquotedblright{}.

\begin{quote}
\emph{fángjiān hěn gānjìng} — The room is clean.
\end{quote}

Its characters come later; for now it is a word for the ear.

\begin{cousinweb}
\emph{gānjìng} is built from gān 'dry' and jìng, both inherited from Middle Chinese, \emph{kan} and \emph{dzjengH} (Guangyun, in Baxter's transcription), glossed 'dry' and 'without dirt'.
\end{cousinweb}

\begin{grammarlens}[title={What you've built}]
An adjective takes a plain \zh{很} where English has \emph{is}: \emph{fángjiān hěn gānjìng}. No \zh{是}.
\end{grammarlens}

\subsection*{Your turn}
\begin{itemize}
\raggedright
  \item \emph{Say it:} \emph{gānjìng}
  \item \emph{Say it:} \emph{gānjìng}, once more
  \item \emph{Say it:} \emph{fángjiān hěn gānjìng} — The room is clean.
  \item \emph{Say it:} \emph{fángjiān bù gānjìng} — The room isn't clean.
  \item \emph{Recall:} say \emph{yǒudiǎnr dà}, then \emph{gānjìng} again
\end{itemize}

\begin{tcolorbox}[breakable,colback=teal!4,colframe=teal!35!black,title={Before you move on}]
What does \emph{gānjìng} mean? (\textquotedblleft{}Clean\textquotedblright{}.) What goes in front of it, \zh{是} or \zh{很}? (\textbf{\zh{很}}.)
\end{tcolorbox}
\section[zāng]{zāng — dirty}

\label{lesson:ZH-C93-zang}



\begin{quote}
Before the new one: last lesson's word, then two things that get dirty.

\begin{itemize}
\raggedright
  \item \emph{Recall:} say \emph{gānjìng}, clean — \textbf{R1}, one lesson back
  \item \emph{Recall:} say \emph{xié}, shoes — \textbf{R3}, durable
  \item \emph{Recall:} say \emph{yīfu}, clothes — \textbf{R3}
\end{itemize}
\end{quote}

\subsection*{You'll want to know: zāng}
\emph{zāng} — \textquotedblleft{}dirty\textquotedblright{}.

\begin{quote}
\emph{xié tài zāng le} — The shoes are too dirty.
\end{quote}

\begin{quote}
\emph{yīfu yǒudiǎnr zāng} — The clothes are a bit dirty.
\end{quote}

Its character comes later; for now it is a word for the ear.

\begin{cousinweb}
The older character for \emph{zāng} is in the Guangyun as \emph{tsangX} (Guangyun, in Baxter's transcription), but only inside a two-character word for a sturdy body, and that reading would give zǎng, not zāng. No dictionary we use records the word 'dirty' in an older form, so the history of \emph{zāng} is not settled.
\end{cousinweb}

\begin{grammarlens}[title={What you've built}]
\emph{gānjìng} and \emph{zāng}, clean and dirty, and three strengths for each: \zh{很}, \zh{有点儿}, \zh{太} … \zh{了}.
\end{grammarlens}

\subsection*{Your turn}
\begin{itemize}
\raggedright
  \item \emph{Say it:} \emph{zāng}
  \item \emph{Say it:} \emph{zāng}, once more
  \item \emph{Say it:} \emph{xié tài zāng le} — The shoes are too dirty.
  \item \emph{Say it:} \emph{yīfu yǒudiǎnr zāng} — The clothes are a bit dirty.
  \item \emph{Recall:} say \emph{gānjìng}, then \emph{zāng} again
\end{itemize}

\begin{tcolorbox}[breakable,colback=teal!4,colframe=teal!35!black,title={Before you move on}]
What does \emph{zāng} mean? (\textquotedblleft{}Dirty\textquotedblright{}.) Say the shoes are too dirty. (\emph{xié tài zāng le}.)
\end{tcolorbox}
\section[Practice]{\zh{太大了}}

\label{lesson:ZH-C93-practice}



\begin{quote}
The chapter closes here, with two people shopping and not finding it.

\begin{itemize}
\raggedright
  \item \emph{Recall:} say \emph{zāng}, dirty — \textbf{R1}, one lesson back
  \item \emph{Recall:} say \emph{zěnmeyàng}, how is it — \textbf{R4}, recognition at distance
  \item \emph{Recall:} say \emph{piányi}, cheap — \textbf{R4}
\end{itemize}
\end{quote}

\subsection*{The exchange}
\begin{quote}
\textbf{A:} \emph{zhège zěnmeyàng} — How's this one?
\end{quote}

\begin{quote}
\textbf{B:} \zh{太大了}
\end{quote}

\begin{quote}
\emph{tài dà le} — Too big.
\end{quote}

\begin{quote}
\textbf{A:} \zh{那个呢}
\end{quote}

\begin{quote}
\emph{nàge ne} — And that one?
\end{quote}

\begin{quote}
\textbf{B:} \zh{有点儿小}
\end{quote}

\begin{quote}
\emph{yǒudiǎnr xiǎo} — A bit small.
\end{quote}

\begin{quote}
\textbf{A:} \zh{这个很好看}
\end{quote}

\begin{quote}
\emph{zhège hěn hǎokàn} — This one's nice.
\end{quote}

\begin{quote}
\textbf{B:} \emph{yǒudiǎnr zāng} — It's a bit dirty.
\end{quote}

\begin{quote}
\textbf{A:} \zh{我们去吃饭}
\end{quote}

\begin{quote}
\emph{wǒmen qù chī fàn} — Let's go and eat, then.
\end{quote}

\begin{quote}
\textbf{B:} \zh{太好了}
\end{quote}

\begin{quote}
\emph{tài hǎo le} — Great!
\end{quote}

Three strengths of complaint, \zh{有点儿}, \zh{很} and \zh{太} … \zh{了}, and one delight to finish: \zh{太好了}.

\begin{tcolorbox}[breakable,skin=enhanced,colback=orange!6,colframe=orange!45!black,boxrule=0.5pt,arc=1mm,left=6pt,right=6pt,top=4pt,bottom=4pt,fonttitle=\bfseries,title={Read this}]
\emph{Read it:} the exchange through once without stopping, then answer the question below
\end{tcolorbox}

\begin{grammarlens}[title={What you've built — three strengths}]
\begin{itemize}
\raggedright
  \item \textbf{\zh{很大}} \emph{hěn dà} — big; a plain statement.
  \item \textbf{\zh{有点儿大}} \emph{yǒudiǎnr dà} — a bit big, and I'd rather it wasn't.
  \item \textbf{\zh{太大了}} \emph{tài dà le} — too big.
  \item \textbf{\zh{太好了}} \emph{tài hǎo le} — great!
\end{itemize}

\textbf{Listen.} \emph{Fángjiān zěnmeyàng?}

\textbf{Speak.} Answer it three ways: \emph{gānjìng}, a bit dirty, too dirty.

\emph{Write it:} \zh{太} — \zh{大}, then the dot
\end{grammarlens}

\subsection*{Your turn}
\begin{itemize}
\raggedright
  \item \emph{Say it:} \zh{太大了} — \emph{tài dà le} — Too big.
  \item \emph{Say it:} \zh{有点儿小} — \emph{yǒudiǎnr xiǎo} — A bit small.
  \item \emph{Say it:} \emph{fángjiān hěn gānjìng} — The room is clean.
  \item \emph{Say it:} the whole exchange, both parts
  \item \emph{Write it:} \zh{太好了}
\end{itemize}

\begin{tcolorbox}[breakable,colback=teal!4,colframe=teal!35!black,title={Before you move on}]
Say it is too big. (\textbf{\zh{太大了}}.) Say it is a bit small. (\textbf{\zh{有点儿小}}.) Say great. (\textbf{\zh{太好了}}.) Say the room is clean. (\emph{fángjiān hěn gānjìng}.)
\end{tcolorbox}
